\section{Branch-Price-and-Cut Algorithm}\label{sec:solution}

Both formulations are solved by branch-price-and-cut. Their internal columns,
internal pricing problem, and machine-side branching rules are identical. They
differ only in the representation of outsourcing: formulation
\textup{(SP2)} keeps the outsourcing decisions in the master, whereas
\textup{(SP1)} invokes a second pricing problem for outsourcing columns.
This section first describes the common node-processing framework and then presents
the two pricing problems and the formulation-specific branching decisions.

\subsection{Restricted master problem and node processing}

At a branch-and-bound node $u$, let $\mathcal R_u\subset\mathcal R$ denote the
current set of internal columns. For \textup{(SP1)}, let
$\mathcal O_u\subset\mathcal O$ denote the current set of outsourcing columns. The
node is processed by repeatedly solving the corresponding restricted master
problem (RMP), generating negative-reduced-cost columns, and reoptimizing the RMP.
In the principal continuous-time configuration, the internal pricing engines are
called in the following order:
\begin{enumerate}[label=(\roman*)]
    \item a tabu-search pricing heuristic;
    \item exact continuous-time bidirectional labeling with ng-route decremental
    state-space relaxation (ng-DSSR).
\end{enumerate}
The time-indexed method is mounted as an alternative exact engine and does not call
the external tabu-search engine first; its inexpensive bucket search, when active,
is internal to the time-indexed pricing procedure.
For \textup{(SP1)}, exact outsourcing pricing is then performed with the
same RMP dual solution. Thus, the pricing loop is closed only after all column
families have been priced exactly.

Let $z_u^{\mathrm{RMP}}$ be the optimal RMP value, let
$\bar c_u^{\mathrm{in}}$ be a certified lower bound on the reduced cost of any
internal column, and, for \textup{(SP1)}, let
$\bar c_u^{\mathrm{out}}$ be the corresponding bound for outsourcing columns. Let
$m_u^{\max}$ denote the machine-count upper bound that remains valid at node $u$. A
valid node lower bound is
\begin{equation}
    \underline z_u
    =
    z_u^{\mathrm{RMP}}
    +
    m_u^{\max}\min\{0,\bar c_u^{\mathrm{in}}\}
    +
    \mathbf 1_{\{\mathrm{SP1}\}}
    \min\{0,\bar c_u^{\mathrm{out}}\}.
    \label{eq:node-dual-bound}
\end{equation}
The machine multiplier in \eqref{eq:node-dual-bound} follows from the node-specific
version of \eqref{eq:machine-capacity}; at most one outsourcing column can be
selected in \textup{(SP1)}. The bound is used only when the relevant pricing
certificates are complete. In particular, a certificate for internal columns alone
cannot close a node of \textup{(SP1)}.

\paragraph{Initial columns and incumbent.}
An initial feasible schedule is obtained by a multi-start adaptive
large-neighborhood search. Its machine sequences are inserted into the global
column pool, together with singleton columns needed to initialize the RMP. When a
root time-indexed preprocessing run is enabled, a small number of elementary
columns with favorable reduced costs are also transferred to the continuous-time
root RMP. At a child node, the restricted column set retains the positive parent
columns and selected low-reduced-cost columns, so that the branch decision is
evaluated from the parent's actual LP support rather than from an unrelated column
sample.

\paragraph{Column identity and cost evaluation.}
The global pool identifies an internal column by its ordered job sequence. The cost
stored for that sequence is the minimum timing cost over its valid time domain, as
defined in \eqref{eq:internal-column-cost}. Bidirectional labeling may recover the same
sequence through different split points. A sequence accepted for the master is
therefore evaluated once over its complete valid domain before it is inserted or
used to update an existing pool entry. This separates route discovery from final
sequence costing and avoids making the master coefficient depend on the particular
bidirectional split used to recover the sequence.

\subsection{Functional representation of internal pricing}
\label{subsec:functional-pricing}

For a fixed partial sequence, its best reduced cost depends on the completion time
of its last job. We represent this relation by a piecewise-linear function (PWLF).
Consider a forward partial sequence ending at job $j$. Its label contains:
\begin{itemize}
    \item the terminal job and a predecessor pointer for sequence recovery;
    \item an ng-memory set and a set of jobs that remain extendable;
    \item a PWLF giving the minimum reduced cost for each feasible completion time
    of $j$.
\end{itemize}
Waiting is handled by the directional closure of the PWLF: a forward function
stores the best value attainable no later than a given completion time, whereas a
backward function stores the symmetric suffix information. Extending a label
through arc $(i,j)$ shifts its time domain by $s_{ij}+p_j$, adds the penalty
function of $j$ and all applicable dual contributions, intersects the result with
the current hard time domain of $j$, and normalizes the resulting function.

The functional representation is important for two reasons. First, it handles
continuous completion times directly and does not create one state for every time
point in the scheduling horizon. Its size is governed by the generated labels and
PWLF segments rather than by an explicit time grid. Second, dominance compares the
complete time--cost trade-off of two partial sequences rather than a single scalar
cost evaluated at an arbitrarily selected time.

\subsection{Heuristic internal pricing}

The heuristic pricing engine is a tabu search over complete machine sequences. Its
seeds are drawn from active RMP columns, prioritizing favorable reduced costs and
positive basic columns. Starting from each seed, the search evaluates three move
families: removing a job, adding a job, and exchanging an included job with an
excluded job. Each iteration applies the best admissible move; a tabu move is
allowed when it improves the best value reached from the current seed.

Sequence costs are updated by functional prefix--suffix composition. Before a move
is accepted, arcs prohibited by the current branching decisions and pricing-only
fixing information are checked directly. Negative-reduced-cost sequences reached
along the accepted search trajectory are collected, deduplicated by sequence, and
returned in increasing reduced-cost order. The heuristic does not provide an
optimality certificate. Failure to find a column therefore transfers control to
the exact internal pricing engine.

\subsection{Exact continuous-time pricing by bidirectional ng-DSSR}
\label{subsec:ng-dssr}

The principal continuous-time configuration studied here combines bidirectional PWLF
labeling, ng-route relaxation, and DSSR. It is exact at termination: intermediate
rounds may admit non-elementary walks, but every returned master column is
elementary, and the ng sets are enlarged whenever a negative non-elementary witness
remains. The ng-memory and decremental-refinement logic is adapted from the route
pricing framework of \citet{martinelli2014pricing}; the state, resource, and cost
representations below are specific to the scheduling pricing problem.

\subsubsection{Bidirectional labeling and midpoint selection}

Forward and backward labels are generated on complementary half-domains separated
by a midpoint $T^{\mathrm{mid}}$. The initial midpoint is derived from the effective
time horizon. A bounded probe evaluates candidate midpoints by running limited
forward and backward expansions. If one direction is substantially heavier, the
probe moves the midpoint toward that direction and, after the imbalance changes
sign, refines the bracket. The selected candidate's labels are reused by the full
labeling run, so the probe is not a separate discarded computation.

Within a DSSR sequence, the previous round supplies feedback on the observed
forward/backward workload. A sufficiently unbalanced round shifts the next probe
seed toward the heavier side; otherwise, the current midpoint is retained as the
next starting point. This mechanism addresses the main failure mode of a fixed
geometric midpoint: long or heterogeneous time windows can make one half-domain
orders of magnitude harder than the other even when their numeric widths are
similar.

\subsubsection{Ng-memory and state-space refinement}

Let $N_j\subseteq\mathcal J\setminus\{j\}$ be the ng-neighborhood of job $j$. For a
forward label $L$ ending at $j$, the ng-memory $M(L)$ contains recently visited
jobs whose repetition remains relevant under the current neighborhoods. Extending
$L$ to job $k$ is forbidden when $k\in M(L)$, and the memory is updated by
intersecting the inherited memory with $N_k$ and adding the newly visited
information required by the ng-route rule. Jobs that cannot be reached again under
the effective time windows are excluded when the initial neighborhoods are
constructed.

The initial neighborhood size grows with the instance size and is restricted to
jobs that can genuinely participate in a repeated visit. After a pricing round,
negative elementary sequences are returned to the master. Negative non-elementary
sequences are retained as DSSR witnesses. For each selected witness, repeated
subsequences identify missing ng pairs; adding these pairs prevents the same
violation from surviving unchanged in the next round. Witnesses are processed in
reduced-cost order, while redundant witnesses already eliminated by earlier
updates are skipped. The algorithm terminates when it finds elementary negative
columns or when an exact round proves that no negative relaxed route remains. The
latter condition also proves that no negative elementary column exists.

\subsubsection{Source-aware incremental dominance}

Our dominance structure adapts the set-dominance graph of
\citet{luo2017dominance}. In that graph, labels with the same earliest service
time and reachable set occupy one node, and the node stores their pointwise lower
envelope over the delivered-quantity resource. Dominance inherited through the
graph can consequently be established by a set of labels even when no individual
label dominates the candidate over its complete domain.

For the present scheduling problem, a separate graph is maintained for each
terminal job and search direction. A node $u$ is identified by a reachable-set key
$R_u$; its complement contains jobs retained by ng-memory and jobs that cannot be
extended from the current functional time domain. The scalar time state used in
the original graph is not a separate key here because feasible completion times
and their reduced costs are represented jointly by the label PWLF. The keys form
a Hasse graph under set inclusion. If $\mathcal A_u$ is the set of active local
labels and $\operatorname{pred}(u)$ is the set of immediate predecessor nodes, the
single envelope maintained at $u$ is
\begin{equation}
    g_u(t)=
    \min\left\{
        \min_{L\in\mathcal A_u} f_L(t),
        \min_{v\in\operatorname{pred}(u)} g_v(t)
    \right\},
    \label{eq:sourced-dominance-envelope}
\end{equation}
where $f_L$ is the frontier of label $L$. Every segment of $g_u$ records either
the local label that supplies it or that the segment is inherited from a
predecessor.

Insertion is incremental. The current $g_u$ first rejects a label that is already
dominated. Otherwise, $f_L$ is merged into $g_u$ by one segment sweep. The merge
returns both the updated envelope and a sparse delta containing only intervals on
which its numerical value decreases. This delta, rather than the complete
envelope, is propagated to immediate successors. Propagation stops at a node when
the delta produces no numerical change. A local label is removed when the merge
shows that its source no longer appears on any segment of the node envelope. An
empty node is deleted and its predecessor--successor links are reconnected.
Under partial dominance, the surviving source intervals are recorded during the
same merge and are applied lazily immediately before the label is expanded or
joined.

This source bookkeeping is the main implementation difference from the
pseudocode of \citet{luo2017dominance}. Their label-removal procedure explicitly
iterates over every label stored at each affected successor node and compares its
function with the updated predecessor envelope. Here, the segment merge itself
identifies which local sources lose their last contributing interval. Dominance
propagation therefore does not scan the node's historical label list or rebuild
separate local, predecessor, and combined envelopes. It still performs
reachable-set graph searches and segment-level envelope scans, so the claim is a
more economical realization of the same set-dominance principle for PWLF
scheduling labels, not a general complexity comparison with the routing
algorithm.

\subsubsection{Completion bounds}

Forward and backward completion-bound functions are computed from a relaxation
that omits ng-memory restrictions while retaining time feasibility and admissible
arcs. They bound the best possible continuation of a partial label. A label is
discarded when its frontier combined with the appropriate completion function
cannot beat the current pricing cutoff.

The bound construction propagates only intervals whose function values have
improved. Forward and backward fixed points are computed first; the auxiliary
transition functions needed by the two labeling directions are then rebuilt once
from the converged functions. This avoids repeatedly merging unchanged portions of
intermediate PWLFs. The same completion information supports reduced-cost arc
fixing after a node is closed. Such fixing uses only node-valid hard domains and
arc restrictions; a dual-profitable window used to accelerate one pricing call is
not written back as permanent subtree fixing evidence.

\subsubsection{Joining forward and backward labels}

Labels are grouped by terminal job and true ng-memory set. Before individual label
pairs are considered, the lower envelopes of a forward and backward group are
joined. If this optimistic group-pair bound is nonnegative, every underlying label
pair is skipped. Otherwise, the algorithm reverts to label-level joining, which
preserves exactness.

For an individual pair, inexpensive compatibility tests precede PWLF composition.
They check the connecting arc, ng-memory intersection, and whether the two partial
sequences already share a job. Visit information is stored in packed machine words,
so these tests usually avoid sequence reconstruction. Only surviving pairs undergo
functional joining. The minimum of the shifted forward frontier plus the backward
frontier is computed by a simultaneous segment scan without constructing a
temporary summed PWLF. Recovered elementary sequences are evaluated over their
complete valid domain before insertion into the master, as explained above.

\subsection{Exact time-indexed pricing}
\label{subsec:time-indexed-pricing}

For integer-time instances, the alternative exact pricing method constructs an
acyclic graph with state $(j,t)$, where $j$ is the last processed job and $t$ is its
completion time. A processing arc appends a job and advances time by its setup and
processing duration; a waiting arc advances time without changing the sequence.
Arc costs include the relevant job, machine, branching, and cut dual terms. A
shortest source--sink path therefore yields a minimum-reduced-cost pseudo-schedule
in the current graph.

The graph can contain repeated job visits. Consequently, negative end states are
ranked and converted to job sequences, but only elementary sequences are returned
as master columns. If the best pseudo-schedule is already nonnegative, the graph
provides a certificate for the entire internal column family. When a
dual-profitable time window is used to shrink a pricing graph, accepted elementary
sequences are evaluated on their full valid domain before their objective
coefficients are stored.

The time-indexed graph also supports reduced-cost fixing. Forward and backward
shortest-path values identify time arcs that cannot belong to an improving column.
A cleanup pass removes waiting or processing states that no longer lead to a useful
extension, and the resulting safe state windows can be inherited by descendants.
The time-indexed method is often effective when the horizon is moderate and the
windows are tight, but its graph size depends directly on the numeric horizon.
This dependence motivates the continuous-time ng-DSSR method.

\paragraph{Rank-1 cuts.}
For the explicit-outsourcing formulation, the time-indexed implementation can be
combined with limited-memory subset-row rank-1 cuts. The cut state is carried in
the pricing labels, and cuts inherited from the parent remain active at child
nodes. This variant strengthens the master relaxation but enlarges the pricing
state space. The source-aware continuous-time ng-DSSR configuration reported here
does not carry active rank-1 cut states: incorporating cut memories into its
dominance keys would destroy the compact envelope structure that provides its main
computational advantage. A separate partial-list continuous-time variant and the
time-indexed rank-1 variant support cut-aware pricing.

\subsection{Outsourcing-column pricing}
\label{subsec:outsourcing-pricing}

This pricing problem is required only by \textup{(SP1)}. For any total
baseline amount $q$, define
\begin{equation}
    V(q)=
    \max\left\{
        \sum_{j\in O}\pi_j:
        O\subseteq\mathcal J,
        \sum_{j\in O}b_j=q
    \right\}.
\end{equation}
At a node with outsourcing-membership branch rows, let $\eta_j$ denote the
corresponding dual contribution for including job $j$ in the outsourcing set and
define $\widetilde\pi_j=\pi_j+\eta_j$; otherwise, let $\eta_j=0$. The definition of
$V(q)$ then uses $\widetilde\pi_j$ in place of $\pi_j$. The pricing objective in
\eqref{eq:outsourcing-subset-pricing} becomes
$\max_q\{V(q)-G(q)\}$. The implementation processes the jobs in a fixed order and
constructs a Pareto frontier of states $(q,v)$, where $q$ is accumulated baseline
outsourcing cost and $v$ is accumulated dual benefit. Excluding job $j$ keeps
$(q,v)$, while including it produces
$(q+b_j,v+\widetilde\pi_j)$. At the same decision stage, state $(q_1,v_1)$
dominates $(q_2,v_2)$ when $q_1\le q_2$ and $v_1\ge v_2$. For every common future
extension, the first state retains at least as much dual benefit and cannot incur a
higher payment because $G$ is nondecreasing. Concavity and piecewise linearity are
therefore unnecessary for this dominance rule; the pricing logic needs only an
evaluator for $G$ at complete states. The implementation used in this paper
evaluates the piecewise-linear tariff defined in
\eqref{eq:outsourcing-tariff-segment}, but replacing that evaluator by any
evaluable nondecreasing aggregate function leaves the dynamic program unchanged.
Branching decisions that require or forbid outsourcing membership are enforced
during the state transitions.

Negative-reduced-cost terminal states are sorted by reduced cost, and the best
configured number of their job sets are added to the outsourcing pool. The same
dynamic program supplies the exact outsourcing reduced-cost certificate used in
\eqref{eq:node-dual-bound}. If outsourcing columns use a more general
inclusion-monotone set cost $C^{\mathrm{out}}(O)$, this dynamic program is replaced
by an exact pricing oracle for that cost; the master formulation is unchanged.

\subsection{Cuts and branching}
\label{subsec:cuts-branching}

The branch-price-and-cut tree uses formulation-specific branching priorities. For
\textup{(SP2)}, the algorithm first branches on the active tariff segment,
then on the number of selected internal-machine columns, and finally on internal
adjacency or directed-arc decisions. For \textup{(SP1)}, tariff branching
is absent. The algorithm branches on the machine count and machine-side
adjacencies/arcs before branching on whether a job belongs to the selected
outsourcing set.

Machine-side branch decisions are enforced both in the RMP and in internal
pricing. An arc-forbidden branch removes the corresponding transition; an
arc-required branch excludes competing transitions that would violate the required
predecessor--successor relation. Outsourcing-membership branches are enforced in
the outsourcing dynamic program and, when a job is required to be outsourced, in
internal pricing as well.

\subsection{Strong branching and Phase-I repair}
\label{subsec:strong-branching}

Strong branching evaluates a bounded set of fractional candidates. For each trial
child, the restricted master is initialized from the parent's positive support and
selected low-reduced-cost columns. If this RMP is infeasible, the implementation
solves a pure Phase-I column-generation problem. Legal columns receive objective
coefficient zero, whereas artificial cover slacks and branch-implied competing
columns receive coefficient one. Heuristic and exact pricing are then performed
with the Phase-I duals. A zero Phase-I objective yields a feasible restricted
master for the child; a positive objective together with complete pricing
certificates proves infeasibility.

After a successful repair, all original column coefficients are restored and the
RMP is reoptimized before the strong-branching score is recorded. This prevents
large artificial penalties from contaminating the pricing duals and separates
feasibility restoration from the child's true objective value. The current
implementation uses one strong-branching screening phase; an additional heuristic
second phase is omitted because its repeated pricing cost was not compensated by
a consistent improvement in candidate selection.

\subsection{Termination and correctness}

At every node, heuristic pricing may add columns but cannot certify closure.
Closure requires the exact internal pricing certificate and, for \textup{(SP1)},
the exact outsourcing certificate. The relaxed ng-route
space contains every elementary internal sequence, so nonnegativity of its exact
minimum reduced cost is sufficient. Time-indexed pricing supplies the corresponding
certificate from the complete acyclic graph. The branch-and-bound node is fathomed
when it is infeasible, when its certified lower bound is no better than the
incumbent, or when the closed RMP solution satisfies all integrality conditions.
Consequently, the algorithm terminates with an optimal solution to the selected
set-covering formulation.
